\documentclass[instructions.tex]{subfiles}
\begin{document}
 
\chapter{Vaines excuses de ceux qui refusent de pardonner.}
 
\primo Si c'était un autre, je lui pardonnerais volontiers ; mais c'est un misérable qui ne mérite pas mon amitié. Et vous qui, après des milliers de péchés, méritez d'être écrasé par toutes les créatures, êtes-vous moins misérable ? S'il ne mérite pas que vous lui pardonniez, Jésus-Christ le mérite ; c'est Jésus-Christ lui-même qui vous demande le pardon pour cet ennemi. Si vous refusez ce pardon à un Dieu, vous n'êtes plus un homme, mais un monstre indigne de vivre.
 
\secundo Je ne puis pardonner ni étouffer le ressentiment, il m'a fait trop de tort. Vous prononcez donc l'arrêt de votre condamnation, lorsque vous dites à Dieu : Pardonnez-nous comme nous pardonnons. Vous vous trompez en disant que vous ne pouvez étouffer le ressentiment. Il est vrai que vous ne pouvez empêcher le ressentiment ; mais vous pouvez y renoncer, n'y pas donner votre consentement. Et comme les pensées qui ne sont pas volontaires ne blessent point la chasteté, ainsi les ressentimens auxquels on ne consent pas, ne blessent point la charité, et n'empêchent pas qu'on ne pardonne.
 
Votre ennemi vous a fait tort ! mais le pardon que vous lui refusez, réparc-t-il le tort qu'il vous a fait ? Loin de-là, votre rancune vous fait à vous-même plus de tort et de mal que la malice de vos ennemis ne peut vous en faire. Un calviniste, ayant reçu sur une image du Sauveur, le trait qu'on lui avait fait tomber mort. Votre prochain est l'image de Dieu : lorsque vous lancez contre lui les traits de votre haine, ce sont autant de traits qui rejaillissent sur votre âme, et qui donnent la mort.
 
Que, et qui oseriez-vous d'un homme à qui on aurait pris une obole, si, pour se venger, il jetait le reste de son bien dans la mer, ou si, pour avoir reçu une égratignure, il se plongeait un poignard dans le sein ? Vous faites à peu près de même, lorsque, pour une perte ou pour une parole, vous donnez, par la rancune, la mort à votre âme. Insensés ! n'estimez-vous votre âme que la valeur d'une somme dont on vous a fait tort, et la mettez-vous en comparaison avec une parole qu'on a dite de vous ?
 
Votre ennemi vous a fait tort ; mais il s'est fait à lui-même plus de tort qu'à vous ; il est plus digne de votre compassion que de votre colere. Si, en vous faisant tort, il tombait dans un précipice et se brisait le corps, vous en auriez compassion. Hélas ! il se porte un coup bien plus funeste, puisqu'en vous persécutant, il précipite son âme dans l'enfer. Voilà ce qui doit vous affliger plus que le tort qu'il vous a fait ; c'est de voir qu'à votre occasion il perd son âme.
 
S'il perd son âme, direz-vous, ce n'est pas ma faute. Ah ! chrétien, où est votre charité et votre foi ? Quoi ! que cette âme se perde sans votre faute, en est-elle moins perdue ? Un Dieu a donné son sang pour la sauver, il a pleuré sa perte, et vous y êtes insensible ! Vous n'êtes pas un homme, mais un démon. Saint François de Sales disait à ce sujet : Je ne sais comment Dieu m'a fait le cœur ; mais s'il m'avait ordonné de haïr un ennemi, je ne pourrais qu'obéir et qu'il ne m'aime point, et qu'il me persécute ? Sans doute, vous y êtes obligé\footnote{Dico vobis: Diligite inimicos vestros. Matth.5.}. Quelle apparence y aurait-il, dit Jésus-Christ, si vous n'aimiez que vos bienfaiteurs et vos amis ?
 
Vous n'êtes pas obligé d'aimer un ennemi d'un amour de confiance, comme vous aimez vos meilleurs amis ; mais vous devez au moins l'aimer d'un amour de patience, souffrir ses défauts, les excuser, n'en point parler pour le diffamer, et ne pas ajouter foi à tout ce qu'on vous en dit. Vous devez l'aimer d'un amour de bienveillance, lui souhaiter du bien, prier pour lui, le servir dans l'occasion, le saluer, lui rendre service dans l'occasion, cela, mais à qui est-ce de commencer ? On raisonne en bon chrétien, c'est à tous les deux, mais chacun prétend que l'autre doit le premier faire les démarches ; mais sans vous embarrasser du devoir d'autrui, faites toujours le vôtre : sans cela, craignez que Dieu ne rejette vos prières et vos sacrifices. Si vous faites une offrande à l'autel, dit Jésus-Christ, et si vous souvenez que votre frère a quelque chose contre vous, laissez là votre offrande, et allez premièrement vous réconcilier. Faut-il, pour un vain point d'honneur, vous exposer à laisser votre âme dans un étal de mort ?
 
Mais je suis au-dessus de lui. Eh qu'importe, il n'est pas ici question de votre rang ; devant Dieu vous êtes égal à votre frère, peut-être encore plus misérable que lui. D'ailleurs, Jésus-Christ n'est-il pas infiniment au-dessus de vous ? Cependant il vous recherche le premier\footnote{Prior dilexit nos, Joan.} N'était-il pas au-dessus de Judas ? cependant il l'embrassa. L'empereur Auguste n'était-il pas au-dessus du traître Cinna ? néanmoins il descendit de son trône pour embrasser ce perfide sénateur, qui en voulait à sa vie.
 
Vous n'êtes pas, à la vérité, obligé de faire des choses indignes de votre rang pour vous reconcilier avec un inférieur ; mais vous êtes obligé de faire les démarches que l'édification publique, que la prudence et la religion exigent de vous.
 
Si, après avoir pris conseil d'un pasteur, selon la conscience, vous demandez une satisfaction d'un tort ou d'une injure atroce, que ce soit pour une fin légitime, dans un esprit de charité, et non pas dans un esprit d'orgueil et de vengeance.
 
\section{Résolutions}
 
\primo Puisque ni la malice de mon ennemi, ni son ingratitude à mon égard, ni la grandeur du tort qu'il m'a fait, ni la disposition où il est encore de me nuire, ne peuvent me dispenser devant Dieu de l'aimer, je renonce à jamais à toute excuse contraire à ce devoir. -- J'aimerai donc cet ennemi ; je prierai pour lui, et je me tiendrai continuellement dans la disposition de l'obliger. 

\prayer{Mon Dieu, donnez-moi la force dont j'ai besoin pour mettre en pratique ces résolutions.}
 
\end{document}
